% ============================================================================
% Chapter IV -- Implementation and Results
% ============================================================================
\chapter{Implementation and Results}
\label{chap:implementation}

\section*{Introduction}

This chapter presents the concrete realization through screenshots, deployment
evidence, and live search demonstrations on the production instance at
\url{https://leakslens.vercel.app}.

% ----------------------------------------------------------------------------
\section{Development Environment}
\label{sec:devenv}
% ----------------------------------------------------------------------------

\begin{table}[H]
\centering
\caption{Development and production environment.}
\label{tab:devenv}
\small
\begin{tabularx}{\textwidth}{@{} l l X @{}}
\toprule
\textbf{Component} & \textbf{Version / Ref.} & \textbf{Usage} \\
\midrule
Development OS    & Windows 11              & Local development and testing \\
Python            & 3.12 (Alpine)           & Backend runtime \\
Node.js           & 20 LTS                  & Frontend build and SSR \\
Docker Desktop    & 27+                     & Local containerization \\
Cursor IDE        & Latest                  & Code editor with AI assistance \\
Git               & 2.x                     & Version control \\
VPS OS            & Ubuntu 22.04 LTS        & Production server \\
VPS Specs         & 91.148.168.36           & Dedicated server \\
Docker (VPS)      & 24+                     & Production containers \\
Nginx Proxy Mgr   & Latest                  & Reverse proxy with SSL \\
Vercel            & Free tier               & Frontend hosting \\
\bottomrule
\end{tabularx}
\end{table}

% ----------------------------------------------------------------------------
\section{User Interface}
\label{sec:ui}
% ----------------------------------------------------------------------------

The frontend uses Next.js~15, Tailwind CSS, and shadcn/ui with a dark theme
suited to extended research sessions.

\subsection{Landing Page (Guest View)}

The landing page presents \leaklens{} to unauthenticated visitors. It
displays the platform's purpose, key features (billions of records, full
credentials, 7 search types, enrichment), and call-to-action buttons for
signing in or creating an account.

\begin{figure}[H]
\centering
\includegraphics[width=\linewidth,keepaspectratio]{landing-guest.png}
\caption{Landing page of \leaklens{} (\url{https://leakslens.vercel.app}).}
\label{fig:landing}
\end{figure}

\subsection{Registration and Login}

Account creation requires email, full name, password, and a valid invite code
when the registration gate is active. The login form accepts email and password;
after successful authentication, the user is redirected to the search
interface. \figref{fig:register} and \figref{fig:login} show both forms on
the same page layout.

\begin{figure}[H]
\centering
\begin{minipage}[t]{0.47\linewidth}
  \centering
  \includegraphics[width=\linewidth,keepaspectratio]{register.png}\\[4pt]
  \captionof{figure}{Registration form with invite-code gate.}
  \label{fig:register}
\end{minipage}\hfill
\begin{minipage}[t]{0.47\linewidth}
  \centering
  \includegraphics[width=\linewidth,keepaspectratio]{login.png}\\[4pt]
  \captionof{figure}{Login form.}
  \label{fig:login}
\end{minipage}
\end{figure}

\subsection{Home Page (Authenticated View)}

Once logged in, the user sees a streamlined home page with the search bar
prominently displayed. Seven search-type chips (Email, Username, IP,
Password, Name, Hash, Domain) allow quick type selection, and a wildcard
toggle enables pattern matching.

\begin{figure}[H]
\centering
\includegraphics[width=\linewidth,keepaspectratio]{home-auth.png}
\caption{Home page with search interface (authenticated user).}
\label{fig:home}
\end{figure}

\subsection{Search Results}

Search results are displayed as cards grouped by source database. Each card
shows the database name, severity badge (critical/high/medium/low), and
credential fields. Sensitive fields (passwords, hashes) are blurred by
default and can be revealed with a click. A copy button is available for each
field.

\begin{figure}[H]
\centering
\includegraphics[width=\linewidth,keepaspectratio]{results.png}
\caption{Search results page showing leak records with severity badges.}
\label{fig:results}
\end{figure}

\subsection{Enrichment Panel}

When searching by email or domain, supplementary intelligence is displayed
in a dedicated enrichment panel: subdomain discovery, WHOIS data, and email
reputation scores.

\begin{figure}[H]
\centering
\includegraphics[width=\linewidth,keepaspectratio]{enrichment.png}
\caption{Enrichment panel showing supplementary intelligence.}
\label{fig:enrichment}
\end{figure}

\subsection{Search History}

The history page lists the user's past searches with the query, type,
result count, and timestamp. A ``Search again'' button allows re-running
any past query.

\begin{figure}[H]
\centering
\includegraphics[width=\linewidth,keepaspectratio]{search-history.png}
\caption{Search history page.}
\label{fig:history}
\end{figure}

\FloatBarrier

% ----------------------------------------------------------------------------
\section{Docker Deployment}
\label{sec:docker-deploy}
% ----------------------------------------------------------------------------

The backend is deployed on the VPS using Docker Compose. The following
screenshots document the build process, running containers, and API health
verification.

\subsection{Docker Compose Build}

Production deployment is started with \texttt{docker compose up -d --build}.
The multi-stage Alpine image compiles Python dependencies in a builder stage,
ships a minimal runtime, and starts the FastAPI and PostgreSQL services with
health-check dependencies.

\begin{figure}[H]
\centering
\includegraphics[width=\linewidth,keepaspectratio]{docker-build.png}
\caption{Successful \texttt{docker compose up -d --build} on the VPS.}
\label{fig:docker-build}
\end{figure}

\subsection{Container Status}

Once the stack is running, both containers report a healthy state. PostgreSQL
remains on the internal Docker network; only the FastAPI service is exposed
to the host on port~8500.

\begin{figure}[H]
\centering
\includegraphics[width=\linewidth,keepaspectratio]{docker-ps.png}
\caption{\texttt{docker ps} output showing \texttt{leaklens-backend-1} and
\texttt{leaklens-postgres-1} in healthy state.}
\label{fig:docker-ps}
\end{figure}

\subsection{Health Check}

The backend exposes a lightweight \texttt{/health} endpoint used by Docker
health checks and for manual verification after deployment.

\begin{figure}[H]
\centering
\includegraphics[width=\linewidth,keepaspectratio]{health-check.png}
\caption{Health check response confirming the API is operational.}
\label{fig:health-check}
\end{figure}

\FloatBarrier

% ----------------------------------------------------------------------------
\section{Search Demonstration}
\label{sec:search-demo}
% ----------------------------------------------------------------------------

This section presents two live queries executed against the production
instance. Results illustrate how \leaklens{} surfaces breached credentials and
organizational exposure from Snusbase and enrichment sources.

\subsection{Password Search}

To demonstrate password-type lookup, the query \texttt{'qwerty@'} was submitted. The platform returned
\textbf{748 records} spanning multiple source databases. Each result card
lists the associated email or username, the exposed password field (blurred by
default in the UI), the originating breach compilation, and a severity badge
--- most entries are rated \emph{critical} because plaintext credentials are
present. \figref{fig:password-search} shows a representative subset of these
results; the full result set confirms how widely reused weak passwords such as
\texttt{qwerty} propagate across unrelated services.

\begin{figure}[H]
\centering
\includegraphics[width=\linewidth,keepaspectratio]{password-search.png}
\caption{Password search for \texttt{qwerty}: 748 leaked credential records
across multiple breach databases.}
\label{fig:password-search}
\end{figure}

\subsection{Domain Search}

A domain search for \texttt{'itbs.tn'} was run to assess exposure tied to an
educational institution. The Search Engine  returned
\textbf{16 breached student records} linked to email addresses under that
domain, together with subdomain and WHOIS context in the enrichment panel.
This type of query helps defenders identify whether institutional accounts
appear in public leak compilations and prioritize password resets or MFA
enforcement for affected users. \figref{fig:search-domain} captures the
domain results and  returned .

\begin{figure}[H]
\centering
\includegraphics[width=\linewidth,keepaspectratio]{domain-search.png}
\caption{Domain search for \texttt{itbs.tn}: 16 breached student records with
subdomain and WHOIS enrichment.}
\label{fig:search-domain}
\end{figure}

\subsection{Name Search}

An individual or employee can search by full name to discover where their
personal data has appeared in breach compilations. \figref{fig:name-search}
shows a name search performed with the author's own name: the platform returned
matching records exposing an associated email address and password in plaintext,
illustrating how quickly personal exposure can be identified.

\begin{figure}[H]
\centering
\includegraphics[width=\linewidth,keepaspectratio]{name-search.png}
\caption{Name search revealing leaked email and password associated with the
queried identity.}
\label{fig:name-search}
\end{figure}

\FloatBarrier

% ----------------------------------------------------------------------------
\section*{Conclusion}

This chapter demonstrated the operational \leaklens{} platform through its
user interface, Docker deployment on the VPS, and three live breach searches.
The password query for \texttt{qwerty@} returned 748 credential records across
multiple databases; the domain query for \texttt{itbs.tn} exposed 16 student
accounts; and a name search confirmed personal exposure by surfacing a linked
email and plaintext password. Together, these results confirm that the
end-to-end search pipeline described in earlier chapters works as designed in
production.
